\documentclass{jlreq}

% font settings
\usepackage[hiragino-pro]{luatexja-preset}
\usepackage{newpxtext, newpxmath}

% math settings
\usepackage{amsmath}

% table settings
\usepackage{tabularray}


% title
\title{RSA/Path/Channel ILP定式化}
\author{森 勇太}
\date{\today}

% document
\begin{document}

\maketitle

\section{モデルの入力}

\subsection{集合・定数}

\begin{tabular}{ll}
  symbol                  & meaning \\
  \hline
  $ E $                   & リンクの集合 \\
  $ S $                   & 周波数スロットの集合 \\
  $ D $                   & 通信要求の集合 \\
  $ M $                   & 変調方式の集合 \\
  $ P^d $                 & 通信要求 $ d \in D $ に対するパスの集合 \\
  $ n^d_p $               & 通信要求 $ d \in D $ に $ p \in P^d $ 割り当てる際の必要なスロット数 \\
  $ C^d_p $               & 通信要求 $ d \in D $ に $ p \in P^d $ 割り当てる際の候補チャネルの集合 \\
  $ \delta^d_{e, p} $     & 通信要求 $ d \in D $ の候補パス $ p \in P^d $ にリンク $ e \in E $ が含まれる場合1 \\
  $ \gamma^d_{s, c, p} $  & 通信要求 $ d \in D $ の候補パス $ p \in P^d $、チャネル $ c \in C^d_p $ にスライス $ s \in S $ が含まれる場合1 \\
\end{tabular}

\subsection{実験パラメータ}

\begin{tabular}{llc}
  name                & type      & related parameter \\
  \hline
  network\_name       & str       & $ E $ \\
  graph               & nx.Graph  & $ E $ \\
  num\_slots          & int       & $ S $ \\
  num\_demands        & int       & $ D $ \\
  demands\_population & list      & $ D $ \\
  k                   & int       & $ P^d $ \\
  path\_algo\_name    & str       & $ P^d $ \\
  alpha:              & float     & $ P^d $ \\
  TimeLimit:          & int       & - \\
  W:                  & dict      & $ n^d_p $ \\
  traffic\_bpsk:      & float     & $ n^d_p $ \\
\end{tabular}

\section{下限値の計算モデル}

\subsection{変数}



\subsection{定式化}

\section{上限値の計算モデル}

\subsection{変数}

\subsection{定式化}

\section{Path/ChannelのILPモデル}

\subsection{変数}

\subsection{定式化}

\end{document}
